\chapter{Conclusiones}
\label{sec:conclusiones}

\section{Desempeño General}
\label{sec:desempenio-general}

El sprint de calidad de SpeechNotes se ejecutó conforme al plan establecido, con una cobertura del 100\% de los requisitos funcionales y una tasa de aprobación del 92.86\% en los casos de prueba ejecutados. Los resultados indican que el sistema se encuentra en un estado funcionalmente estable, con oportunidades de mejora identificadas en areas de manejo de errores y experiencia de usuario.

\section{Hallazgos Clave}
\label{sec:hallazgos-clave}

\begin{enumerate}
    \item \textbf{Funcionalidad Core Robusta:} Los requisitos funcionales relacionados con autenticación, transcripción en vivo, upload de archivos, consulta y edición de transcripciones, búsqueda y formateo con IA pasaron todas las pruebas.

    \item \textbf{Problemas de UX Detectados:} Se identificaron 3 bugs de severidad Minor relacionados con la experiencia de usuario (tooltip duplicado, icono de play/pausa inconsistente, placeholder en registro), los cuales no afectan la funcionalidad pero impactan la calidad percibida del producto.

    \item \textbf{Mejora Significativa en Calidad de Código:} La integración de SonarCloud permitió identificar y corregir code smells. La configuración del pipeline CI/CD asegura que la calidad del código se mantenga en el tiempo.

    \item \textbf{Automatización como Inversión:} Las suites de pytest y Cypress configuradas proporcionan una base sólida para regression testing en sprints futuros.

    \item \textbf{Bugs Identificados:} Se encontraron 8 bugs (3 corregidos, 5 pendientes). Los bugs críticos y mayores fueron atendidos durante el sprint, mientras que los menores quedaron registrados para atención futura.
\end{enumerate}

\section{Recomendaciones}
\label{sec:recomendaciones}

\begin{enumerate}
    \item \textbf{Corregir bugs pendientes en el c\'odigo:} Los 5 bugs abiertos (BUG-002, BUG-005, BUG-006, BUG-008, BUG-009) deben ser registrados como tickets en Jira y atendidos en el siguiente sprint de desarrollo.
    \item \textbf{Pruebas de carga:} Realizar pruebas de carga formal con k6 o Locust para validar el comportamiento bajo estrés.
    \item \textbf{Documentación de API:} Completar la documentación de los endpoints REST de la API.
    \item \textbf{Pruebas de seguridad:} Integrar OWASP ZAP en el pipeline CI/CD.
    \item \textbf{Monitoreo de SonarCloud:} Mantener el quality gate activo en el pipeline de CI/CD para prevenir regresiones.
    \item \textbf{Ampliar automatización:} Incrementar la cobertura de pruebas automatizadas, especialmente en flujos de error y casos límite.
    \item \textbf{Mejoras de UX:} Implementar tooltip contextual en lugar del modal de error de micrófono.
\end{enumerate}

\section{Lecciones Aprendidas}
\label{sec:lecciones}

\begin{enumerate}
    \item La planificación detallada (matriz de rastreabilidad, casos de prueba) antes de la ejecución redujo significativamente la ambigüedad durante las pruebas.
    \item La integración temprana de SonarCloud permitió detectar code smells antes de que se acumularan.
    \item El uso de Jira para el seguimiento de bugs facilitó la comunicación entre QA y desarrollo.
    \item Para sprints futuros, se recomienda asignar tiempo adicional para pruebas exploratorias y de accesibilidad.
\end{enumerate}
